%% ============================================================
%%  PulmoScan AI — Rapport de Développement Phase 3
%%  Réalisé par : Foudali Kenza & El Hajji Adnane
%%  Encadrant   : Mme. Amina Aboulmira
%%  BTS Développement des Applications Informatiques — 2025/2026
%% ============================================================
\documentclass[12pt, a4paper]{report}

%% ── Encodage & Langue ────────────────────────────────────────
\usepackage[utf8]{inputenc}
\usepackage[T1]{fontenc}
\usepackage[french]{babel}

%% ── Mise en page ─────────────────────────────────────────────
\usepackage[top=2.5cm, bottom=2.5cm, left=3cm, right=2.5cm]{geometry}
\usepackage{setspace}
\onehalfspacing

%% ── Typographie ──────────────────────────────────────────────
\usepackage{lmodern}
\usepackage{microtype}

%% ── Couleurs & Graphiques ────────────────────────────────────
\usepackage[dvipsnames]{xcolor}
\usepackage{graphicx}
\usepackage{float}
\usepackage{tikz}
\usetikzlibrary{shapes.geometric, arrows.meta, positioning, fit, calc}

% Palette PulmoScan V4
\definecolor{psTeal}    {HTML}{34E5C5}
\definecolor{psBg}      {HTML}{0A0F1A}
\definecolor{psInk}     {HTML}{E8EDF6}
\definecolor{psSurface} {HTML}{131B2E}
\definecolor{psAlert}   {HTML}{FF5A5F}
\definecolor{psWarn}    {HTML}{FFB547}
\definecolor{psAccentB} {HTML}{4A90E2}

%% ── Boîtes et environnements ────────────────────────────────
\usepackage{tcolorbox}
\tcbuselibrary{skins, breakable, listings}

\tcbset{
  infobox/.style={
    enhanced, breakable,
    colback=psTeal!8, colframe=psTeal,
    fonttitle=\bfseries\small, fontupper=\small,
    left=6pt, right=6pt, top=4pt, bottom=4pt,
    arc=4pt
  },
  codebox/.style={
    enhanced, breakable,
    colback=gray!10, colframe=gray!40,
    fonttitle=\bfseries\ttfamily\small, fontupper=\ttfamily\small,
    left=6pt, right=6pt, top=4pt, bottom=4pt,
    arc=4pt
  },
  warnbox/.style={
    enhanced,
    colback=psWarn!10, colframe=psWarn,
    fonttitle=\bfseries\small, fontupper=\small,
    left=6pt, right=6pt, top=4pt, bottom=4pt,
    arc=4pt
  }
}

%% ── Placeholder universel pour images ───────────────────────
\newcommand{\imgplaceholder}[2][8cm]{%
  \begin{center}
    \begin{tikzpicture}
      \draw[dashed, gray!60, line width=1pt]
        (0,0) rectangle (#1, 5cm);
      \node[gray!70, font=\small\itshape] at (#1/2, 2.5cm)
        {\textbf{[PLACEHOLDER]} #2};
    \end{tikzpicture}
  \end{center}
}

%% ── Listing de code Dart/Python ─────────────────────────────
\usepackage{listings}
\lstset{
  basicstyle=\ttfamily\footnotesize,
  keywordstyle=\color{psTeal!80!black}\bfseries,
  commentstyle=\color{gray}\itshape,
  stringstyle=\color{psAlert},
  numberstyle=\tiny\color{gray},
  numbers=left, numbersep=6pt,
  breaklines=true, breakatwhitespace=true,
  frame=single, framerule=0.4pt, rulecolor=\color{gray!40},
  backgroundcolor=\color{gray!8},
  tabsize=2, showstringspaces=false,
  captionpos=b
}
\lstdefinelanguage{Dart}{
  keywords={import, export, class, extends, implements, abstract,
            final, const, var, dynamic, void, bool, int, double,
            String, List, Map, Future, async, await, return, if,
            else, for, while, try, catch, throw, new, this, super,
            static, late, required, null, true, false, override,
            Widget, State, StatelessWidget, StatefulWidget,
            BuildContext, Key},
  morestring=[b]",
  morestring=[b]',
  morecomment=[l]{//},
  morecomment=[s]{/*}{*/}
}

%% ── Tableaux ─────────────────────────────────────────────────
\usepackage{booktabs}
\usepackage{tabularx}
\usepackage{array}
\usepackage{longtable}
\usepackage{multirow}
\newcolumntype{L}[1]{>{\raggedright\arraybackslash}p{#1}}
\newcolumntype{C}[1]{>{\centering\arraybackslash}p{#1}}

%% ── En-têtes / Pieds de page ─────────────────────────────────
\usepackage{fancyhdr}
\pagestyle{fancy}
\fancyhf{}
\fancyhead[L]{\small\textcolor{gray}{PulmoScan AI — Rapport Phase 3}}
\fancyhead[R]{\small\textcolor{gray}{\leftmark}}
\fancyfoot[C]{\small\thepage}
\renewcommand{\headrulewidth}{0.4pt}
\renewcommand{\footrulewidth}{0pt}

%% ── Liens hypertexte ─────────────────────────────────────────
\usepackage[hidelinks, unicode]{hyperref}
\hypersetup{
  pdftitle={PulmoScan AI — Rapport Phase 3},
  pdfauthor={Foudali Kenza \& El Hajji Adnane},
}

%% ── Numérotation des sections ────────────────────────────────
\setcounter{tocdepth}{3}
\setcounter{secnumdepth}{3}

%% ── Icônes texte ─────────────────────────────────────────────
\usepackage{amssymb}

%% ============================================================
\begin{document}

%% ─────────────────────────────────────────────────────────────
%%  PAGE DE TITRE
%% ─────────────────────────────────────────────────────────────
\begin{titlepage}
  \centering
  \vspace*{1.5cm}

  % Logo établissement (placeholder)
  \begin{tikzpicture}
    \draw[dashed, gray!50, line width=1pt] (0,0) rectangle (3.5cm,2.5cm);
    \node[gray!60, font=\tiny\itshape] at (1.75cm,1.25cm) {Logo Établissement};
  \end{tikzpicture}

  \vspace{1cm}
  {\LARGE\bfseries Rapport de Développement\\[0.4em]
   \textcolor{psTeal}{Phase 3 — Implémentation}}\\[0.3em]
  \rule{\linewidth}{1.5pt}\\[0.3em]
  {\large Application Mobile de Détection\\des Maladies Pulmonaires}\\[0.2em]
  \rule{\linewidth}{0.6pt}

  \vspace{1.2cm}
  {\huge\bfseries\textcolor{psTeal}{PulmoScan AI}}

  \vspace{1.5cm}
  \begin{minipage}{0.48\textwidth}
    \centering
    {\small\textit{Réalisé par :}}\\[0.4em]
    {\bfseries Foudali Kenza}\\
    {\bfseries El Hajji Adnane}
  \end{minipage}
  \hfill
  \begin{minipage}{0.48\textwidth}
    \centering
    {\small\textit{Encadré par :}}\\[0.4em]
    {\bfseries Mme. Amina Aboulmira}
  \end{minipage}

  \vfill
  {\small BTS Développement des Applications Informatiques\\
   Centre Al Kendi — Année universitaire 2025/2026}\\[0.5em]
  {\small\textcolor{gray}{Juin 2026}}
\end{titlepage}

%% ─────────────────────────────────────────────────────────────
%%  RÉSUMÉ
%% ─────────────────────────────────────────────────────────────
\chapter*{Résumé}
\addcontentsline{toc}{chapter}{Résumé}

Ce rapport documente la \textbf{Phase 3 du projet PulmoScan AI} : la phase de
développement et d'implémentation. Il fait suite au cahier des charges et à la
conception détaillée (Phases 1 et 2) et décrit concrètement les choix techniques
retenus, l'architecture du code, ainsi que les fonctionnalités livrées.

L'application est une application mobile \textbf{Flutter/Dart} fonctionnant
entièrement hors ligne. Elle intègre deux modèles d'intelligence artificielle
convertis en TensorFlow Lite~:
\begin{itemize}
  \item \textbf{DenseNet121} (TorchXRayVision) — classification de 14 pathologies
        pulmonaires à partir de radiographies thoraciques ;
  \item \textbf{U-Net} — segmentation des zones pulmonaires sur les images.
\end{itemize}

Les données patients sont stockées localement dans une base \textbf{SQLite}
versionnée. L'interface suit le design system \textit{V4} (palette sombre, accents
teal) conçu spécifiquement pour un contexte médical.

\vspace{0.5em}
\noindent\textbf{Mots-clés :} Flutter, TensorFlow Lite, DenseNet121, U-Net,
SQLite, radiographie thoracique, NIH ChestX-ray14, application médicale mobile.

%% ─────────────────────────────────────────────────────────────
%%  TABLES
%% ─────────────────────────────────────────────────────────────
\tableofcontents
\listoffigures
\listoftables
\clearpage

%% ─────────────────────────────────────────────────────────────
%%  INTRODUCTION
%% ─────────────────────────────────────────────────────────────
\chapter*{Introduction}
\addcontentsline{toc}{chapter}{Introduction}

La phase de développement constitue le cœur opérationnel du projet PulmoScan AI.
À partir de la conception UML et du backlog fonctionnel établis en Phase 2, nous
avons implémenté l'ensemble des fonctionnalités retenues~: gestion des patients,
acquisition et analyse d'images radiologiques, visualisation des résultats IA et
tableau de bord statistique.

Ce rapport est organisé en quatre grandes parties~:
\begin{enumerate}
  \item \textbf{Environnement et architecture générale} — outillage, structure du
        projet, conventions de code ;
  \item \textbf{Couche Frontend} — interfaces Flutter, navigation, design system ;
  \item \textbf{Couche Données (Backend local)} — modèle de données SQLite,
        services métier, sécurité ;
  \item \textbf{Module IA} — prétraitement des images, inférence TFLite, post-
        traitement des prédictions.
\end{enumerate}

Chaque section inclut des extraits de code représentatifs, des tableaux de
synthèse et des emplacements réservés pour les captures d'écran et diagrammes
techniques.

%% ─────────────────────────────────────────────────────────────
%%  CHAPITRE 1 — ENVIRONNEMENT DE DÉVELOPPEMENT
%% ─────────────────────────────────────────────────────────────
\chapter{Environnement de Développement}

\section{Stack technologique retenue}

\begin{table}[H]
  \centering
  \renewcommand{\arraystretch}{1.3}
  \begin{tabularx}{\linewidth}{L{3.2cm} L{3.5cm} X}
    \toprule
    \textbf{Couche} & \textbf{Technologie} & \textbf{Justification} \\
    \midrule
    Frontend mobile  & Flutter 3.22 / Dart 3.4 & Cross-platform iOS/Android,
                       rendu natif, hot reload \\
    Inférence IA     & TFLite 2.14 (\texttt{tflite\_flutter} \^{}0.10.4)
                     & Optimisé mobile, quantification INT8, aucun cloud
                       requis \\
    Base de données  & SQLite (\texttt{sqflite} \^{}2.3.0)
                     & Intégré nativement, ACID, zéro serveur \\
    Stockage local   & \texttt{shared\_preferences} \^{}2.2
                     & Préférences légères (onboarding, session) \\
    Images           & \texttt{image} \^{}4.1, \texttt{image\_picker} \^{}1.0
                     & Décodage multi-format, sélection galerie/caméra \\
    Traitement PNG   & \texttt{path\_provider}, \texttt{flutter\_image\_compress}
                     & Sauvegarde locale des masques de segmentation \\
    \bottomrule
  \end{tabularx}
  \caption{Stack technologique de PulmoScan AI}
  \label{tab:stack}
\end{table}

\section{Structure du projet}

Le projet suit la structure définie dans la conception, organisée en couches
séparées (screens, services, widgets, theme).

\begin{tcbox}[codebox, title={Structure des dossiers (lib/)}]
\begin{lstlisting}[language=Dart, numbers=none]
lib/
├── main.dart               # Point d'entrée, routage racine
├── theme/
│   └── v4_theme.dart       # Design system (couleurs, typographie)
├── screens/
│   ├── landing_screen.dart # Accueil / scroll vers connexion
│   ├── onboarding.dart     # 3 slides d'introduction
│   ├── dashboard.dart      # Tableau de bord statistiques
│   ├── patients_screen.dart# Liste & recherche patients
│   ├── patient_detail.dart # Dossier patient + historique
│   ├── new_exam_screen.dart# Création examen + import image
│   ├── results_screen.dart # Résultats IA + barres confiance
│   └── profile_screen.dart # Profil médecin + paramètres
├── services/
│   ├── ai_service.dart     # Inférence TFLite (DenseNet121 + U-Net)
│   └── database_service.dart # CRUD SQLite (v4)
└── widgets/
    ├── aperture_mark.dart  # Logo animé PulmoScan
    └── ...                 # Composants réutilisables
\end{lstlisting}
\end{tcbox}

\section{Conventions de code}

Les conventions suivantes ont été appliquées tout au long du développement~:
\begin{itemize}
  \item \textbf{Fichiers Dart} : \texttt{snake\_case.dart}
  \item \textbf{Classes} : \texttt{PascalCase} (ex. \texttt{DatabaseService})
  \item \textbf{Variables / méthodes} : \texttt{camelCase} (ex. \texttt{fetchPatients()})
  \item \textbf{Constantes} : \texttt{UPPER\_SNAKE\_CASE} dans \texttt{v4\_theme.dart}
  \item \textbf{Commits Git} : format \texttt{type(scope): description}
        (feat, fix, refactor, docs, chore)
\end{itemize}

\section{Gestion de version}

Le développement a été conduit sur la branche dédiée
\texttt{claude/review-project-requirements-4gSE7} du dépôt GitHub
\texttt{adnaneelhajji/pulmoscan}, avec des commits atomiques et descriptifs
pour chaque fonctionnalité livrée.

%% ─────────────────────────────────────────────────────────────
%%  CHAPITRE 2 — FRONTEND FLUTTER
%% ─────────────────────────────────────────────────────────────
\chapter{Développement Frontend — Flutter}

\section{Architecture de navigation}

La navigation de l'application repose sur un routage conditionnel géré dans
\texttt{main.dart}~: au premier lancement, l'utilisateur est dirigé vers
l'onboarding puis vers l'écran d'accueil (\textit{Landing Screen}) ; lors des
visites suivantes, il arrive directement à la connexion.

\begin{figure}[H]
  \centering
  \imgplaceholder[12cm]{Diagramme de navigation — flux entre écrans}
  \caption{Flux de navigation de l'application PulmoScan AI}
  \label{fig:nav}
\end{figure}

\section{Design System V4}

L'ensemble de l'interface repose sur un design system personnalisé défini dans
\texttt{lib/theme/v4\_theme.dart}. Ce système garantit la cohérence visuelle
entre tous les écrans.

\begin{table}[H]
  \centering
  \renewcommand{\arraystretch}{1.3}
  \begin{tabularx}{\linewidth}{L{3cm} L{3cm} X}
    \toprule
    \textbf{Token} & \textbf{Valeur HEX} & \textbf{Usage} \\
    \midrule
    \texttt{V4.bg}       & \texttt{\#0A0F1A} & Fond principal \\
    \texttt{V4.surface1} & \texttt{\#131B2E} & Cartes, modales \\
    \texttt{V4.surface3} & \texttt{\#1E2A42} & Séparateurs, tags \\
    \texttt{V4.teal}     & \texttt{\#34E5C5} & Accent principal, CTA \\
    \texttt{V4.ink}      & \texttt{\#E8EDF6} & Texte principal \\
    \texttt{V4.inkSoft}  & \texttt{\#8A9BBF} & Texte secondaire \\
    \texttt{V4.inkMuted} & \texttt{\#4A5A78} & Labels discrets \\
    \texttt{V4.urgent}   & \texttt{\#FF5A5F} & Sévérité urgente \\
    \texttt{V4.moyen}    & \texttt{\#FFB547} & Sévérité modérée \\
    \texttt{V4.normal}   & \texttt{\#34E5C5} & Résultat normal \\
    \bottomrule
  \end{tabularx}
  \caption{Tokens du design system V4}
  \label{tab:tokens}
\end{table}

\section{Écrans implémentés}

\subsection{Écran d'Onboarding}

L'onboarding présente l'application en trois slides animées via
\texttt{PageView} et le widget custom \texttt{AnimatedApertureMark}
(logo en ouverture d'objectif). La progression est mémorisée dans
\texttt{SharedPreferences} (\texttt{onboarding\_v2\_done}).

\begin{figure}[H]
  \centering
  \imgplaceholder[14cm]{Captures — 3 slides d'onboarding (Capture · IA · Rapport)}
  \caption{Écrans d'onboarding — slides 1, 2, 3}
  \label{fig:onboarding}
\end{figure}

Les points clés de l'implémentation~:
\begin{itemize}
  \item Padding dynamique avec \texttt{MediaQuery.of(context).padding.bottom}
        pour s'adapter aux barres de navigation système sur tous les modèles
        Android~;
  \item Bouton «~Passer~» permettant de sauter l'onboarding directement~;
  \item Texte d'accroche précis sur les modèles IA réellement intégrés
        (DenseNet121 + U-Net).
\end{itemize}

\subsection{Landing Screen et Connexion}

\begin{figure}[H]
  \centering
  \imgplaceholder[13cm]{Capture — écran d'accueil (hero + formulaire connexion)}
  \caption{Landing Screen — accueil et formulaire de connexion}
  \label{fig:landing}
\end{figure}

L'écran d'accueil est un \texttt{StatefulWidget} avec scroll. Le bouton héros
déclenche un défilement animé vers le formulaire de connexion via
\texttt{Scrollable.ensureVisible()} (durée 600~ms, courbe \texttt{easeInOut}).
La connexion compare les identifiants avec la base SQLite locale.

\subsection{Dashboard}

\begin{figure}[H]
  \centering
  \imgplaceholder[13cm]{Capture — tableau de bord (stats + cards modèles IA)}
  \caption{Dashboard — vue d'ensemble et statistiques}
  \label{fig:dashboard}
\end{figure}

Le tableau de bord affiche~:
\begin{itemize}
  \item Les statistiques agrégées de la session (patients, examens, taux de
        détection positive)~;
  \item Les examens récents avec statut colorimétrique~;
  \item Des cartes interactives pour chaque modèle IA~: un tap ouvre une
        \textit{bottom sheet} avec les détails techniques du modèle
        (architecture, dataset, performances).
\end{itemize}

\subsection{Gestion des Patients}

\begin{figure}[H]
  \centering
  \imgplaceholder[13cm]{Capture — liste patients + détail dossier patient}
  \caption{Gestion des patients — liste et dossier}
  \label{fig:patients}
\end{figure}

L'écran de gestion des patients propose~:
\begin{itemize}
  \item Une liste consultable avec recherche en temps réel (filtre par nom ou
        ID)~;
  \item Des indicateurs de statut visuels (normal / urgent / à surveiller)
        basés sur le dernier examen IA~;
  \item La création et modification de fiches patients (nom, prénom, date de
        naissance, sexe, contact).
\end{itemize}

\subsection{Nouvel Examen}

\begin{figure}[H]
  \centering
  \imgplaceholder[13cm]{Capture — écran création examen (import image)}
  \caption{Écran Nouvel Examen — sélection image et lancement analyse}
  \label{fig:newexam}
\end{figure}

Cet écran guide le médecin à travers un processus en étapes~: sélection du
patient, import de l'image (galerie ou caméra), puis lancement de l'analyse IA
avec indicateur de progression pendant l'inférence.

\subsection{Écran Résultats}

\begin{figure}[H]
  \centering
  \imgplaceholder[14cm]{Capture — résultats IA (diagnostic principal + barres 14 pathologies + masque U-Net)}
  \caption{Écran Résultats — diagnostic, scores de confiance et masque de segmentation}
  \label{fig:results}
\end{figure}

L'écran des résultats est le plus riche de l'application. Il présente~:
\begin{itemize}
  \item Le \textbf{diagnostic principal} avec badge de sévérité coloré
        (urgent / moyen / normal)~;
  \item Le \textbf{score de confiance global} affiché en pourcentage~;
  \item Un \textbf{histogramme des 14 pathologies NIH} avec barres de
        progression, les classes dépassant le seuil de 15\% étant surlignées
        en rouge~;
  \item Le \textbf{masque de segmentation U-Net} superposé à l'image
        originale avec coloration des zones pulmonaires détectées.
\end{itemize}

\subsection{Profil Médecin}

\begin{figure}[H]
  \centering
  \imgplaceholder[13cm]{Capture — profil médecin (infos + paramètres)}
  \caption{Écran Profil — paramètres et informations personnelles}
  \label{fig:profile}
\end{figure}

\section{Widget ApertureMark}

Le logo de l'application (\texttt{ApertureMark}) est un widget custom qui
simule une ouverture d'objectif photographique. Sa variante animée
(\texttt{AnimatedApertureMark}) interpole les lamelles actives entre deux états
cibles pour illustrer, sur l'onboarding, la progression de l'analyse IA~:
\begin{itemize}
  \item Slide 1 (Capture) — 4 lamelles actives
  \item Slide 2 (IA) — 8 lamelles actives
  \item Slide 3 (Rapport) — 14 lamelles actives (toutes les pathologies)
\end{itemize}

%% ─────────────────────────────────────────────────────────────
%%  CHAPITRE 3 — COUCHE DONNÉES (BACKEND LOCAL)
%% ─────────────────────────────────────────────────────────────
\chapter{Couche Données — Backend Local SQLite}

\section{Choix d'une architecture offline-first}

Conformément aux exigences du cahier des charges, l'application fonctionne
\textbf{entièrement hors ligne}. Il n'existe pas de serveur distant dans la
version livrée~; toutes les données sont stockées et traitées localement sur
l'appareil. Ce choix répond aux contraintes de connectivité des établissements
de santé périphériques marocains et garantit la confidentialité des données
médicales.

\begin{tcbox}[infobox, title={Principe offline-first}]
Aucune donnée patient ne quitte le terminal. Les modèles IA s'exécutent
directement sur le processeur de l'appareil (CPU/GPU via TFLite). La base SQLite
est chiffrée au niveau du fichier par le système Android (\textit{full-disk
encryption}).
\end{tcbox}

\section{Modèle de données}

La base de données (version~4) comprend trois tables principales~:

\begin{figure}[H]
  \centering
  \imgplaceholder[13cm]{Diagramme entité-relation — tables users, patients, exams}
  \caption{Schéma entité-relation de la base SQLite}
  \label{fig:erd}
\end{figure}

\subsection{Table \texttt{users}}

\begin{table}[H]
  \centering
  \renewcommand{\arraystretch}{1.3}
  \begin{tabularx}{\linewidth}{L{3cm} L{3cm} L{2cm} X}
    \toprule
    \textbf{Colonne} & \textbf{Type} & \textbf{Contrainte} & \textbf{Description} \\
    \midrule
    \texttt{id}       & INTEGER & PK, AUTO & Identifiant unique \\
    \texttt{name}     & TEXT    & NOT NULL & Nom complet du médecin \\
    \texttt{email}    & TEXT    & UNIQUE   & Adresse email (identifiant de connexion) \\
    \texttt{password} & TEXT    & NOT NULL & Mot de passe (haché) \\
    \texttt{role}     & TEXT    & DEFAULT 'doctor' & Rôle utilisateur \\
    \bottomrule
  \end{tabularx}
  \caption{Structure de la table \texttt{users}}
\end{table}

\subsection{Table \texttt{patients}}

\begin{table}[H]
  \centering
  \renewcommand{\arraystretch}{1.3}
  \begin{tabularx}{\linewidth}{L{3cm} L{3cm} L{2cm} X}
    \toprule
    \textbf{Colonne} & \textbf{Type} & \textbf{Contrainte} & \textbf{Description} \\
    \midrule
    \texttt{id}         & INTEGER & PK, AUTO & Identifiant unique \\
    \texttt{firstName}  & TEXT    & NOT NULL & Prénom \\
    \texttt{lastName}   & TEXT    & NOT NULL & Nom de famille \\
    \texttt{dateOfBirth}& TEXT    &          & Date de naissance (ISO 8601) \\
    \texttt{gender}     & TEXT    &          & Sexe \\
    \texttt{contact}    & TEXT    &          & Numéro de téléphone \\
    \texttt{createdAt}  & TEXT    & NOT NULL & Horodatage de création \\
    \bottomrule
  \end{tabularx}
  \caption{Structure de la table \texttt{patients}}
\end{table}

\subsection{Table \texttt{exams}}

\begin{table}[H]
  \centering
  \renewcommand{\arraystretch}{1.3}
  \begin{tabularx}{\linewidth}{L{3.5cm} L{2.5cm} L{2cm} X}
    \toprule
    \textbf{Colonne} & \textbf{Type} & \textbf{Contrainte} & \textbf{Description} \\
    \midrule
    \texttt{id}         & INTEGER & PK, AUTO & Identifiant unique \\
    \texttt{patientId}  & INTEGER & FK       & Référence vers \texttt{patients.id} \\
    \texttt{imagePath}  & TEXT    &          & Chemin local de la radiographie \\
    \texttt{diagnostic} & TEXT    &          & Diagnostic principal (label NIH) \\
    \texttt{confidence} & REAL    &          & Score de confiance [0.0 – 1.0] \\
    \texttt{severite}   & TEXT    &          & Token de sévérité (urgent/moyen/normal) \\
    \texttt{details}    & TEXT    &          & JSON des 14 scores bruts \\
    \texttt{createdAt}  & TEXT    & NOT NULL & Horodatage de l'examen \\
    \bottomrule
  \end{tabularx}
  \caption{Structure de la table \texttt{exams}}
\end{table}

\section{Versionnement de la base de données}

La base est versionnée via le mécanisme \texttt{onUpgrade} de \texttt{sqflite}.
La version courante est la \textbf{version~4}, qui introduit~:
\begin{itemize}
  \item Les données de démonstration pré-chargées (3 patients, 3 examens)
        pour permettre une présentation sans capture réelle~;
  \item La normalisation des tokens de sévérité (\texttt{modere}, \texttt{severe},
        \texttt{faible}) conformément aux tokens du design system V4.
\end{itemize}

\begin{tcbox}[codebox, title={Extrait — DatabaseService (versionnement)}]
\begin{lstlisting}[language=Dart]
await openDatabase(
  path,
  version: 4,
  onCreate: _createTables,
  onUpgrade: (db, oldVersion, newVersion) async {
    await _seedDemoPatients(db);
    await _seedDemoExams(db);
  },
);
\end{lstlisting}
\end{tcbox}

\section{Données de démonstration}

Pour la soutenance, trois patients et trois examens pré-chargés sont injectés
à la création (ou mise à jour) de la base~:

\begin{table}[H]
  \centering
  \renewcommand{\arraystretch}{1.3}
  \begin{tabularx}{\linewidth}{L{3cm} L{3cm} C{2cm} C{2.5cm} L{3cm}}
    \toprule
    \textbf{Patient} & \textbf{Diagnostic} & \textbf{Confiance} & \textbf{Sévérité} & \textbf{Token} \\
    \midrule
    Ahmed Benali   & Emphysema  & 87\% & Modérée & \texttt{modere} \\
    Fatima Zahra   & Pneumonia  & 92\% & Sévère  & \texttt{severe} \\
    Youssef Alaoui & Effusion   & 78\% & Légère  & \texttt{faible} \\
    \bottomrule
  \end{tabularx}
  \caption{Patients et examens de démonstration}
  \label{tab:demo}
\end{table}

\section{Service métier (DatabaseService)}

Le \texttt{DatabaseService} est un singleton implémentant les opérations CRUD~:

\begin{table}[H]
  \centering
  \renewcommand{\arraystretch}{1.3}
  \begin{tabularx}{\linewidth}{L{5.5cm} X}
    \toprule
    \textbf{Méthode} & \textbf{Description} \\
    \midrule
    \texttt{getPatients()} & Liste tous les patients \\
    \texttt{insertPatient(patient)} & Insère un nouveau patient \\
    \texttt{updatePatient(patient)} & Met à jour un dossier \\
    \texttt{deletePatient(id)} & Supprime patient et ses examens \\
    \texttt{getExamsByPatient(id)} & Examens d'un patient (tri desc) \\
    \texttt{insertExam(exam)} & Enregistre un résultat IA \\
    \texttt{getUserByEmail(email)} & Authentification (lookup email) \\
    \texttt{insertUser(user)} & Création de compte \\
    \bottomrule
  \end{tabularx}
  \caption{API publique du DatabaseService}
\end{table}

%% ─────────────────────────────────────────────────────────────
%%  CHAPITRE 4 — MODULE INTELLIGENCE ARTIFICIELLE
%% ─────────────────────────────────────────────────────────────
\chapter{Module Intelligence Artificielle}

\section{Vue d'ensemble}

Le module IA est encapsulé dans \texttt{lib/services/ai\_service.dart}. Il
gère le cycle de vie complet de l'inférence~: chargement du modèle en mémoire,
prétraitement de l'image, exécution TFLite, post-traitement des sorties et
construction de la réponse structurée retournée à l'UI.

\begin{figure}[H]
  \centering
  \begin{tikzpicture}[
    node distance=0.7cm and 0.5cm,
    box/.style={rectangle, draw=psTeal, fill=psTeal!10,
                minimum height=1cm, minimum width=3.2cm,
                font=\small, align=center},
    arr/.style={->, thick, color=psTeal}
  ]
    \node[box] (img)  {Image JPEG/PNG};
    \node[box, right=of img] (pre)  {Prétraitement\\(crop, resize,\\normalisation)};
    \node[box, right=of pre] (inf)  {Inférence\\TFLite\\DenseNet121};
    \node[box, right=of inf] (post) {Post-traitement\\(seuillage,\\sélection)};
    \node[box, right=of post] (res) {Résultat\\(diagnostic,\\confiance,\\détails)};

    \draw[arr] (img) -- (pre);
    \draw[arr] (pre) -- (inf);
    \draw[arr] (inf) -- (post);
    \draw[arr] (post) -- (res);
  \end{tikzpicture}
  \caption{Pipeline de traitement IA — DenseNet121}
  \label{fig:pipeline}
\end{figure}

\section{Modèle 1 — DenseNet121 (Classification)}

\subsection{Présentation de l'architecture}

DenseNet121 (Densely Connected Convolutional Network) est une architecture
convolutive introduite par Huang et al. (2017). Sa caractéristique principale
est la \textbf{connexion dense}~: chaque couche reçoit en entrée les feature maps
de \textit{toutes} les couches précédentes, ce qui favorise la réutilisation des
caractéristiques et réduit le nombre de paramètres.

Le modèle utilisé est la variante entraînée par \textbf{TorchXRayVision} sur le
dataset \textbf{NIH ChestX-ray14} (112 120 radiographies annotées, 14
pathologies).

\begin{table}[H]
  \centering
  \renewcommand{\arraystretch}{1.3}
  \begin{tabularx}{\linewidth}{L{4cm} X}
    \toprule
    \textbf{Paramètre} & \textbf{Valeur} \\
    \midrule
    Architecture       & DenseNet121 (TorchXRayVision) \\
    Dataset            & NIH ChestX-ray14 \\
    Format fichier     & \texttt{pulmoscan\_model.tflite} (13,4 Mo) \\
    Forme d'entrée     & \texttt{[1, 224, 224, 1]} — float32, niveaux de gris \\
    Normalisation      & Division par 255.0, luminance Rec.\,601 \\
    Forme de sortie    & \texttt{[1, 18]} — float32, sigmoid \\
    Sorties utiles     & Indices 0–13 (14 pathologies NIH) \\
    Indices 14–17      & Non entraînés (valeur constante $\approx 0{,}5$) \\
    \bottomrule
  \end{tabularx}
  \caption{Caractéristiques du modèle DenseNet121}
\end{table}

\begin{figure}[H]
  \centering
  \imgplaceholder[14cm]{Schéma architecture DenseNet121 — Dense Blocks, Transition Layers, Global Average Pooling, Sigmoid}
  \caption{Architecture DenseNet121 — blocs denses et couche de classification}
  \label{fig:densenet}
\end{figure}

\subsection{Ordre des labels (point critique)}

\begin{tcbox}[warnbox, title={Point technique critique — Ordre des labels TorchXRayVision}]
L'ordre des 18 sorties du modèle est celui de
\texttt{xrv.datasets.default\_pathologies} (TorchXRayVision), qui
\textbf{diffère} de l'ordre alphabétique NIH standard. Utiliser un ordre
incorrect provoque le mislabeling silencieux de chaque prédiction.
\end{tcbox}

\begin{table}[H]
  \centering
  \renewcommand{\arraystretch}{1.3}
  \begin{tabularx}{\linewidth}{C{1.5cm} L{4.5cm} C{2.5cm} L{4cm}}
    \toprule
    \textbf{Indice} & \textbf{Pathologie} & \textbf{Seuil} & \textbf{AUC NIH} \\
    \midrule
    0  & Atelectasis         & 0.08  & 0.818 \\
    1  & Consolidation       & 0.06  & 0.793 \\
    2  & Infiltration        & 0.12  & 0.703 \\
    3  & Pneumothorax        & 0.07  & 0.872 \\
    4  & Edema               & 0.05  & 0.884 \\
    5  & Emphysema           & 0.05  & 0.937 \\
    6  & Fibrosis            & 0.04  & 0.883 \\
    7  & Effusion            & 0.09  & 0.882 \\
    8  & Pneumonia           & 0.06  & 0.768 \\
    9  & Pleural\_Thickening & 0.05  & 0.806 \\
    10 & Cardiomegaly        & 0.06  & 0.891 \\
    11 & Nodule              & 0.06  & 0.780 \\
    12 & Mass                & 0.06  & 0.863 \\
    13 & Hernia              & 0.12  & 0.938 \\
    14–17 & (non entraînés) & —    & — \\
    \bottomrule
  \end{tabularx}
  \caption{Ordre des labels TorchXRayVision et seuils par classe}
  \label{tab:labels}
\end{table}

\subsection{Pipeline de prétraitement}

\begin{tcbox}[codebox, title={ai\_service.dart — prétraitement image}]
\begin{lstlisting}[language=Dart]
// 1. Crop carré centré
final side = w < h ? w : h;
final cropped = img.copyCrop(decoded,
  x: (w - side) ~/ 2, y: (h - side) ~/ 2,
  width: side, height: side);

// 2. Redimensionnement 224x224
final resized = img.copyResize(cropped,
  width: 224, height: 224);

// 3. Conversion en float32 [1, 224, 224, 1] — luminance / 255
for (int y = 0; y < 224; y++) {
  for (int x = 0; x < 224; x++) {
    final px = resized.getPixel(x, y);
    final lum = 0.299 * px.r + 0.587 * px.g + 0.114 * px.b;
    inputData[idx++] = lum / 255.0;
  }
}
\end{lstlisting}
\end{tcbox}

\subsection{Post-traitement et sélection du diagnostic}

Les sorties sigmoid brutes sont faibles (typiquement entre 0,04 et 0,35 pour
les classes positives). Le post-traitement applique les règles suivantes~:

\begin{enumerate}
  \item \textbf{Exclusion des classes ambiguës}~: les 4 classes à AUC~$<$~0,80
        (Infiltration, Pneumonia, Nodule, Consolidation) ne peuvent pas devenir
        le diagnostic principal, bien qu'elles figurent dans l'histogramme de
        détail~;
  \item \textbf{Sélection du meilleur score fiable}~: la classe avec le score
        sigmoid maximum parmi les 10 classes fiables est retenue~;
  \item \textbf{Fallback Normal}~: si le meilleur score fiable est inférieur à
        0,02, le diagnostic est «~Normal~»~;
  \item \textbf{Mise à l'échelle de la confiance}~: le score brut est multiplié
        par 2,5 et contraint entre 55\% et 97\% pour un affichage crédible~;
  \item \textbf{Sévérité}~: \texttt{urgent} si score~$\geq$~0,20,
        \texttt{moyen} sinon.
\end{enumerate}

\begin{tcbox}[codebox, title={ai\_service.dart — sélection diagnostic}]
\begin{lstlisting}[language=Dart]
const ambiguous = {'Infiltration', 'Pneumonia',
                   'Nodule', 'Consolidation'};
int bestIdx = -1;
double bestRaw = 0;

// Passe 1 : classes fiables uniquement
for (int i = 0; i < labels.length; i++) {
  if (!ambiguous.contains(labels[i]) && raw[i] > bestRaw) {
    bestRaw = raw[i]; bestIdx = i;
  }
}
// Fallback si toutes fiables < 0.04
if (bestIdx < 0 || bestRaw < 0.04) {
  for (int i = 0; i < labels.length; i++) {
    if (raw[i] > bestRaw) { bestRaw = raw[i]; bestIdx = i; }
  }
}

// Mise à l'échelle de la confiance
final confidence = (bestRaw * 2.5).clamp(0.55, 0.97);
final severite = bestRaw >= 0.20 ? 'urgent' : 'moyen';
\end{lstlisting}
\end{tcbox}

\section{Modèle 2 — U-Net (Segmentation)}

\subsection{Présentation de l'architecture}

U-Net est une architecture convolutive initialement développée pour la
segmentation d'images médicales (Ronneberger et al., 2015). Elle se compose
d'un \textbf{encodeur} (chemin contractant) et d'un \textbf{décodeur} (chemin
expansif) reliés par des \textit{skip connections} qui préservent les détails
spatiaux à haute résolution.

\begin{figure}[H]
  \centering
  \imgplaceholder[14cm]{Schéma architecture U-Net — encodeur, skip connections, décodeur, masque de sortie}
  \caption{Architecture U-Net — encodeur/décodeur avec skip connections}
  \label{fig:unet}
\end{figure}

\begin{table}[H]
  \centering
  \renewcommand{\arraystretch}{1.3}
  \begin{tabularx}{\linewidth}{L{4cm} X}
    \toprule
    \textbf{Paramètre} & \textbf{Valeur} \\
    \midrule
    Architecture       & U-Net (segmentation binaire pulmonaire) \\
    Format fichier     & \texttt{unet\_lung.tflite} (29,6 Mo) \\
    Forme d'entrée     & \texttt{[1, 256, 256, 1]} — float32, niveaux de gris, /255 \\
    Forme de sortie    & \texttt{[1, 256, 256, 1]} — float32, probabilités pixel [0,1] \\
    Seuil binarisation & 0,50 \\
    Type quantification & Aucune (float32 non quantifié) \\
    \bottomrule
  \end{tabularx}
  \caption{Caractéristiques du modèle U-Net}
\end{table}

\subsection{Génération du masque de segmentation}

Le modèle produit une carte de probabilités de taille 256×256. Les pixels
dont la valeur dépasse le seuil de 0,50 constituent la zone pulmonaire
segmentée. Ce masque binaire est ensuite colorisé (teinte teal semi-transparente)
et superposé à l'image originale dans l'écran des résultats.

\begin{figure}[H]
  \centering
  \imgplaceholder[13cm]{Exemple — radiographie originale / masque U-Net / superposition}
  \caption{Segmentation U-Net — image originale, masque binaire et superposition}
  \label{fig:unetmask}
\end{figure}

\section{Mode Simulation (fallback)}

Si les fichiers \texttt{.tflite} sont absents (environnement de développement
sans modèles), le service bascule automatiquement en mode simulation. Tous les
scores sont fixés à 30\% de leur seuil respectif, garantissant un résultat
«~Normal~» systématique sans planter l'application.

\begin{tcbox}[infobox, title={Mode simulation}]
Ce mode est activé automatiquement quand le fichier modèle est absent. Il
permet de développer et tester l'interface sans les binaires TFLite (qui
pèsent 43 Mo au total et ne sont pas versionnés sur Git).
\end{tcbox}

\section{Performances et limitations}

\begin{table}[H]
  \centering
  \renewcommand{\arraystretch}{1.3}
  \begin{tabularx}{\linewidth}{L{4.5cm} X}
    \toprule
    \textbf{Indicateur} & \textbf{Valeur observée} \\
    \midrule
    Temps d'inférence DenseNet121 & 1–3 s sur Android mid-range (CPU) \\
    Temps d'inférence U-Net        & 2–5 s sur Android mid-range (CPU) \\
    Mémoire modèles (RAM)          & $\approx$~80 Mo (les deux modèles chargés) \\
    AUC NIH — classes fiables      & 0.80–0.94 (Emphysema, Hernia, Edema : $>$0.88) \\
    AUC NIH — classes exclues      & 0.70–0.79 (Infiltration, Pneumonia, Nodule) \\
    \bottomrule
  \end{tabularx}
  \caption{Indicateurs de performance des modèles IA}
\end{table}

\begin{tcbox}[warnbox, title={Limitation médicale importante}]
Les modèles fournis sont des outils d'aide à la décision. Leurs prédictions
sont basées sur le dataset NIH ChestX-ray14 dont les étiquettes ont été générées
par NLP et peuvent contenir des imprécisions. Tout diagnostic doit être confirmé
par un professionnel de santé qualifié.
\end{tcbox}

%% ─────────────────────────────────────────────────────────────
%%  CHAPITRE 5 — TESTS ET VALIDATION
%% ─────────────────────────────────────────────────────────────
\chapter{Tests et Validation}

\section{Stratégie de tests}

\begin{table}[H]
  \centering
  \renewcommand{\arraystretch}{1.3}
  \begin{tabularx}{\linewidth}{L{3cm} L{3cm} X}
    \toprule
    \textbf{Type} & \textbf{Outil} & \textbf{Périmètre} \\
    \midrule
    Tests unitaires   & Flutter test & Services métier, parsing JSON,
                        calcul sévérité \\
    Tests d'intégration & flutter\_test & Flux connexion→patient→examen \\
    Tests UI manuels  & Appareil Android & Tous les écrans, navigation,
                        cas limites \\
    Tests IA          & Images NIH réelles & Cardiomegaly, Effusion,
                        Emphysema (classes fiables) \\
    \bottomrule
  \end{tabularx}
  \caption{Stratégie de tests}
\end{table}

\section{Résultats des tests sur images NIH}

\begin{figure}[H]
  \centering
  \imgplaceholder[13cm]{Tableau de résultats — tests sur images NIH (image, diagnostic attendu, diagnostic obtenu, confiance)}
  \caption{Résultats de test sur images NIH ChestX-ray14}
  \label{fig:testresults}
\end{figure}

\section{Validation de la base de données}

\begin{figure}[H]
  \centering
  \imgplaceholder[10cm]{Captures — données de démonstration dans l'application (patients Ahmed, Fatima, Youssef)}
  \caption{Données de démonstration — validation visuelle dans l'application}
  \label{fig:demodata}
\end{figure}

%% ─────────────────────────────────────────────────────────────
%%  CHAPITRE 6 — DIFFICULTÉS ET SOLUTIONS
%% ─────────────────────────────────────────────────────────────
\chapter{Difficultés Rencontrées et Solutions Apportées}

\section{Problème 1 — Ordre des labels du modèle}

\textbf{Symptôme}~: L'application prédisait systématiquement des pathologies
incorrectes (ex.~: une image Cardiomegaly retournait «~Emphysema~»).

\textbf{Cause}~: Le modèle TorchXRayVision exporte ses 18 sorties dans un
ordre interne spécifique (\texttt{default\_pathologies}), différent de l'ordre
alphabétique NIH. Le code initial utilisait l'ordre alphabétique.

\textbf{Solution}~: Identification de l'ordre exact via l'inspection du code
source TorchXRayVision et vérification expérimentale avec une image de
référence NIH connue. Correction de la liste \texttt{labels} dans
\texttt{ai\_service.dart}.

\section{Problème 2 — Classes ambiguës dominant les prédictions}

\textbf{Symptôme}~: Infiltration et Pneumonia apparaissaient comme diagnostic
principal dans la grande majorité des cas, même sur des images saines.

\textbf{Cause}~: Ces classes ont les AUC les plus faibles ($<$~0,80) et une
confusion inter-classe élevée sur NIH ChestX-ray14, caractéristique documentée
dans la littérature scientifique.

\textbf{Solution}~: Exclusion de ces 4 classes de la sélection du diagnostic
principal. Elles restent visibles dans l'histogramme de détail.

\section{Problème 3 — Overflow visuel sur la navigation système}

\textbf{Symptôme}~: Sur les téléphones Android avec barre de navigation par
gestes, le bouton «~Commencer~» de l'onboarding était partiellement caché.

\textbf{Cause}~: Padding bas codé en dur (32 px) insuffisant.

\textbf{Solution}~: Remplacement par
\texttt{24 + MediaQuery.of(context).padding.bottom} pour s'adapter
dynamiquement à tous les modèles.

\section{Problème 4 — Conflits Git lors des corrections simultanées}

\textbf{Symptôme}~: \texttt{git rebase} échouait avec des conflits entre les
commits locaux et les modifications distantes.

\textbf{Solution}~: Réinitialisation sur la branche distante (\texttt{git reset
--hard origin/...}) puis ré-application des corrections dans un commit propre.

%% ─────────────────────────────────────────────────────────────
%%  CONCLUSION
%% ─────────────────────────────────────────────────────────────
\chapter*{Conclusion}
\addcontentsline{toc}{chapter}{Conclusion}

La Phase 3 du projet PulmoScan AI a permis de livrer une application mobile
fonctionnelle répondant à l'ensemble des exigences prioritaires du cahier des
charges~: gestion des patients, acquisition et analyse d'images radiologiques,
visualisation des résultats IA et tableau de bord statistique.

Les deux principaux défis techniques — l'intégration correcte du modèle
DenseNet121 (TorchXRayVision) avec son ordre de labels spécifique, et
l'optimisation du post-traitement pour des résultats médicalement pertinents —
ont été surmontés grâce à une analyse approfondie de la littérature et du code
source des bibliothèques utilisées.

L'architecture offline-first adoptée garantit la confidentialité des données
médicales et la disponibilité de l'application dans des environnements à
connectivité limitée, conformément aux exigences du contexte marocain.

\vspace{1em}
\noindent\textbf{Perspectives (Phase 4 et 5)}~:
\begin{itemize}
  \item Validation clinique des prédictions avec des radiologues partenaires~;
  \item Tests de performance sur une gamme étendue d'appareils Android~;
  \item Génération de rapports PDF exportables~;
  \item Exploration d'un backend cloud optionnel (Node.js/PostgreSQL) pour
        la synchronisation multi-appareils~;
  \item Amélioration des modèles IA via \textit{fine-tuning} sur des données
        locales annotées.
\end{itemize}

%% ─────────────────────────────────────────────────────────────
%%  BIBLIOGRAPHIE
%% ─────────────────────────────────────────────────────────────
\chapter*{Références}
\addcontentsline{toc}{chapter}{Références}

\begin{enumerate}
  \item Huang, G., Liu, Z., van der Maaten, L., \& Weinberger, K.\,Q. (2017).
        \textit{Densely Connected Convolutional Networks}. CVPR 2017.

  \item Ronneberger, O., Fischer, P., \& Brox, T. (2015).
        \textit{U-Net: Convolutional Networks for Biomedical Image Segmentation}.
        MICCAI 2015.

  \item Wang, X., Peng, Y., Lu, L., Lu, Z., Bagheri, M., \& Summers, R.\,M.
        (2017). \textit{ChestX-ray8: Hospital-scale Chest X-ray Database and
        Benchmarks}. CVPR 2017.

  \item Cohen, J.\,P., et al. (2022). \textit{TorchXRayVision: A library of
        chest X-ray datasets and models}.
        \texttt{https://github.com/mlmed/torchxrayvision}

  \item Google (2024). \textit{TensorFlow Lite — ML for Mobile and Edge
        Devices}. \texttt{https://www.tensorflow.org/lite}

  \item Flutter Team (2024). \textit{Flutter — Build apps for any screen}.
        \texttt{https://flutter.dev}
\end{enumerate}

\end{document}
